% GENERATED FILE. Edit canonical lessons, then run npm run generate:books.
% canonical-source-hash: fnv1a64:f831597f8702233a
% canonical-lessons: SA-C486-upavasati, SA-C486-pranayamah, SA-C486-dhyanam, SA-C486-abhyangah, SA-C486-brahmamuhurtah

\chapter{Fasting and Meditation}
\label{ch:sa-fasting-and-meditation}
\hlchaptermodality{486}

\begin{chapteropening}
\textbf{By the end of this chapter:} \emph{I can say what I do to stay well.}

Complete the last lesson of chapter 486: I can say what I do to stay well.
\end{chapteropening}

\section[\texorpdfstring{\sk{उपवसति} (upavasati)}{उपवसति (upavasati)}]{\sk{उपवसति} (upavasati) — fasts}

\label{lesson:SA-C486-upavasati}



\begin{quote}
Before the new one: say the Sanskrit for relief, then the Sanskrit for a rejuvenating tonic.
\end{quote}

\subsection*{You'll want to know: \sk{उपवसति}}
\textbf{\sk{उपवसति}} — \emph{upavasati} — \textquotedblleft{}fasts\textquotedblright{}.

\begin{quote}
\sk{रुग्णः} \sk{बालः} \sk{अद्य} \sk{उपवसति।}
\end{quote}

\begin{quote}
\emph{rugṇaḥ bālaḥ adya upavasati.}
\end{quote}

\begin{quote}
The sick child fasts today.
\end{quote}

Past (he …): \sk{सः} \sk{उपोषितवान्}, \emph{saḥ upoṣitavān}.

\begin{cousinweb}
\textbf{\sk{उपवसति}} (\emph{upavasati}) is built from \emph{upa} and the root √vas, \textquotedblleft{}to dwell\textquotedblright{}, the root of \textbf{\sk{वसति}}: to stay in abstinence; MW gives \textquotedblleft{}to abstain from food, fast\textquotedblright{} (MW).
\end{cousinweb}

\begin{grammarlens}[title={What you've built}]
One more word for describing what you lived and felt.
\end{grammarlens}

\subsection*{Your turn}
\begin{itemize}
\raggedright
  \item \emph{Say it:} \emph{upavasati}
  \item \emph{Say it:} \emph{upavasati}, once more
  \item \emph{Say it:} \emph{rugṇaḥ bālaḥ adya upavasati.} — The sick child fasts today.
  \item \emph{Recall:} say the Sanskrit for relief, then the Sanskrit for a rejuvenating tonic, then say \emph{upavasati} again
\end{itemize}

\begin{tcolorbox}[breakable,colback=teal!4,colframe=teal!35!black,title={Before you move on}]
What does \sk{उपवसति} mean? (\textquotedblleft{}Fasts\textquotedblright{}.) Say it once more.
\end{tcolorbox}
\section[\texorpdfstring{\sk{प्राणायामः} (prāṇāyāmaḥ)}{प्राणायामः (prāṇāyāmaḥ)}]{\sk{प्राणायामः} (prāṇāyāmaḥ) — breath control}

\label{lesson:SA-C486-pranayamah}



\begin{quote}
Before the new one: say the Sanskrit for a rejuvenating tonic, then the Sanskrit for fasts.
\end{quote}

\subsection*{You'll want to know: \sk{प्राणायामः}}
\textbf{\sk{प्राणायामः}} — \emph{prāṇāyāmaḥ} — \textquotedblleft{}breath control\textquotedblright{}.

\begin{quote}
\sk{पिता} \sk{प्रतिदिनं} \sk{प्राणायामं} \sk{करोति।}
\end{quote}

\begin{quote}
\emph{pitā pratidinaṃ prāṇāyāmaṃ karoti.}
\end{quote}

\begin{quote}
Father does breath control every day.
\end{quote}

A masculine noun.

\begin{cousinweb}
\textbf{\sk{प्राणायामः}} (\emph{prāṇāyāmaḥ}) is built from \emph{prāṇa}, \textquotedblleft{}breath\textquotedblright{}, and \emph{āyāma}, \textquotedblleft{}restraint\textquotedblright{}, from \emph{ā} and the root √yam; MW names it as the \textquotedblleft{}breath-exercises\textquotedblright{} of the Saṃdhyā (MW).
\end{cousinweb}

\begin{grammarlens}[title={What you've built}]
One more word for describing what you lived and felt.
\end{grammarlens}

\subsection*{Your turn}
\begin{itemize}
\raggedright
  \item \emph{Say it:} \emph{prāṇāyāmaḥ}
  \item \emph{Say it:} \emph{prāṇāyāmaḥ}, once more
  \item \emph{Say it:} \emph{pitā pratidinaṃ prāṇāyāmaṃ karoti.} — Father does breath control every day.
  \item \emph{Recall:} say the Sanskrit for a rejuvenating tonic, then the Sanskrit for fasts, then say \emph{prāṇāyāmaḥ} again
\end{itemize}

\begin{tcolorbox}[breakable,colback=teal!4,colframe=teal!35!black,title={Before you move on}]
What does \sk{प्राणायामः} mean? (\textquotedblleft{}Breath control\textquotedblright{}.) Say it once more.
\end{tcolorbox}
\section[\texorpdfstring{\sk{ध्यानम्} (dhyānam)}{ध्यानम् (dhyānam)}]{\sk{ध्यानम्} (dhyānam) — meditation}

\label{lesson:SA-C486-dhyanam}



\begin{quote}
Before the new one: say the Sanskrit for fasts, then the Sanskrit for breath control.
\end{quote}

\subsection*{You'll want to know: \sk{ध्यानम्}}
\textbf{\sk{ध्यानम्}} — \emph{dhyānam} — \textquotedblleft{}meditation\textquotedblright{}.

\begin{quote}
\sk{ध्यानेन} \sk{मनः} \sk{शान्तं} \sk{भवति।}
\end{quote}

\begin{quote}
\emph{dhyānena manaḥ śāntaṃ bhavati.}
\end{quote}

\begin{quote}
Meditation calms the mind.
\end{quote}

A neuter noun.

\begin{cousinweb}
\textbf{\sk{ध्यानम्}} (\emph{dhyānam}) comes from the root √dhyai, \textquotedblleft{}to think of\textquotedblright{}, the root of \textbf{\sk{ध्यायति}}; MW gives \textquotedblleft{}meditation, thought, reflection\textquotedblright{} (MW).
\end{cousinweb}

\begin{grammarlens}[title={What you've built}]
One more word for describing what you lived and felt.
\end{grammarlens}

\subsection*{Your turn}
\begin{itemize}
\raggedright
  \item \emph{Say it:} \emph{dhyānam}
  \item \emph{Say it:} \emph{dhyānam}, once more
  \item \emph{Say it:} \emph{dhyānena manaḥ śāntaṃ bhavati.} — Meditation calms the mind.
  \item \emph{Recall:} say the Sanskrit for fasts, then the Sanskrit for breath control, then say \emph{dhyānam} again
\end{itemize}

\begin{tcolorbox}[breakable,colback=teal!4,colframe=teal!35!black,title={Before you move on}]
What does \sk{ध्यानम्} mean? (\textquotedblleft{}Meditation\textquotedblright{}.) Say it once more.
\end{tcolorbox}
\section[\texorpdfstring{\sk{अभ्यंगः} (abhyaṅgaḥ)}{अभ्यंगः (abhyaṅgaḥ)}]{\sk{अभ्यंगः} (abhyaṅgaḥ) — an oil massage}

\label{lesson:SA-C486-abhyangah}



\begin{quote}
Before the new one: say the Sanskrit for breath control, then the Sanskrit for meditation.
\end{quote}

\subsection*{You'll want to know: \sk{अभ्यंगः}}
\textbf{\sk{अभ्यंगः}} — \emph{abhyaṅgaḥ} — \textquotedblleft{}an oil massage\textquotedblright{}.

\begin{quote}
\sk{शीते} \sk{तैलस्य} \sk{अभ्यंगः} \sk{सुखं} \sk{ददाति।}
\end{quote}

\begin{quote}
\emph{śīte tailasya abhyaṅgaḥ sukhaṃ dadāti.}
\end{quote}

\begin{quote}
In the cold an oil massage brings comfort.
\end{quote}

A masculine noun.

\begin{cousinweb}
\textbf{\sk{अभ्यंगः}} (\emph{abhyaṅgaḥ}) is built from \emph{abhi} and the root √añj, \textquotedblleft{}to anoint\textquotedblright{}; MW gives \textquotedblleft{}rubbing with unctuous substances, inunction\textquotedblright{} (MW).
\end{cousinweb}

\begin{grammarlens}[title={What you've built}]
One more word for describing what you lived and felt.
\end{grammarlens}

\subsection*{Your turn}
\begin{itemize}
\raggedright
  \item \emph{Say it:} \emph{abhyaṅgaḥ}
  \item \emph{Say it:} \emph{abhyaṅgaḥ}, once more
  \item \emph{Say it:} \emph{śīte tailasya abhyaṅgaḥ sukhaṃ dadāti.} — In the cold an oil massage brings comfort.
  \item \emph{Recall:} say the Sanskrit for breath control, then the Sanskrit for meditation, then say \emph{abhyaṅgaḥ} again
\end{itemize}

\begin{tcolorbox}[breakable,colback=teal!4,colframe=teal!35!black,title={Before you move on}]
What does \sk{अभ्यंगः} mean? (\textquotedblleft{}An oil massage\textquotedblright{}.) Say it once more.
\end{tcolorbox}
\section[\texorpdfstring{\sk{ब्रह्ममुहूर्तः} (brahmamuhūrtaḥ)}{ब्रह्ममुहूर्तः (brahmamuhūrtaḥ)}]{\sk{ब्रह्ममुहूर्तः} (brahmamuhūrtaḥ) — the hour before dawn}

\label{lesson:SA-C486-brahmamuhurtah}



\begin{quote}
Before the new one: say the Sanskrit for meditation, then the Sanskrit for an oil massage.
\end{quote}

\subsection*{You'll want to know: \sk{ब्रह्ममुहूर्तः}}
\textbf{\sk{ब्रह्ममुहूर्तः}} — \emph{brahmamuhūrtaḥ} — \textquotedblleft{}the hour before dawn\textquotedblright{}.

\begin{quote}
\sk{पितामहः} \sk{ब्रह्ममुहूर्ते} \sk{उत्तिष्ठति।}
\end{quote}

\begin{quote}
\emph{pitāmahaḥ brahmamuhūrte uttiṣṭhati.}
\end{quote}

\begin{quote}
Grandfather gets up before dawn.
\end{quote}

A masculine noun.

\begin{cousinweb}
\textbf{\sk{ब्रह्ममुहूर्तः}} (\emph{brahmamuhūrtaḥ}) is built from \emph{brahman} and \textbf{\sk{मुहूर्तः}}, \textquotedblleft{}an hour\textquotedblright{}; MW gives only \textquotedblleft{}a particular hour of the day\textquotedblright{}; \textquotedblleft{}the hour before dawn\textquotedblright{} is the sense it has in modern use (MW).
\end{cousinweb}

\begin{grammarlens}[title={What you've built}]
One more word for describing what you lived and felt.
\end{grammarlens}

\subsection*{Your turn}
\begin{itemize}
\raggedright
  \item \emph{Say it:} \emph{brahmamuhūrtaḥ}
  \item \emph{Say it:} \emph{brahmamuhūrtaḥ}, once more
  \item \emph{Say it:} \emph{pitāmahaḥ brahmamuhūrte uttiṣṭhati.} — Grandfather gets up before dawn.
  \item \emph{Recall:} say the Sanskrit for meditation, then the Sanskrit for an oil massage, then say \emph{brahmamuhūrtaḥ} again
\end{itemize}

\begin{tcolorbox}[breakable,colback=teal!4,colframe=teal!35!black,title={Before you move on}]
What does \sk{ब्रह्ममुहूर्तः} mean? (\textquotedblleft{}The hour before dawn\textquotedblright{}.) Say it once more.
\end{tcolorbox}
